% AUTHORED, NOT EXTRACTED: the v1.1 manuscript contains no propositional
% Hilbert calculus.  Canonical source: Mendelson, Introduction to Mathematical
% Logic (1997), axiom schemas A1-A3 with modus ponens (after Lukasiewicz);
% cf. Kleene, Introduction to Metamathematics (1952).  The negation-primitive
% third schema is chosen to match the book's connective conventions, in which
% negation is primitive in the active classical language.
\calculussheet{sheet:hilbert-cpl}{Hilbert Calculus for Classical Propositional Logic}{Language: the $\neg/\to$ fragment of $\mathsf{Fm}_{\mathsf{CPL}}$ (metavariables $\varphi,\psi,\chi$ range over formulas of that fragment; the book's remaining primitives $\bot,\land,\lor$ are reached through the translation of Definition~\ref{def:p02-hilbert-translation} when this calculus is compared with the full language); sequent shape: no sequents --- Hilbert-style derivability $\Gamma\vdash\varphi$ via finite sequences of formulas, each a member of $\Gamma$, an instance of an axiom schema, or obtained from earlier lines by modus ponens; negation: primitive (the third schema governs it directly); equality: none (propositional language).}
\hbrule{rule:hilbert-cpl-a1}{A1}{}{\varphi\to(\psi\to\varphi)}{\none}
\hbrule{rule:hilbert-cpl-a2}{A2}{}{(\varphi\to(\psi\to\chi))\to((\varphi\to\psi)\to(\varphi\to\chi))}{\none}
\hbrule{rule:hilbert-cpl-a3}{A3}{}{(\neg\psi\to\neg\varphi)\to((\neg\psi\to\varphi)\to\psi)}{\none}
\hbrule{rule:hilbert-cpl-mp}{MP}{\varphi\to\psi\qquad\varphi}{\psi}{\none}

\paragraph{Deduction theorem.}
For every set \(\Gamma\) and all formulas \(\varphi,\psi\) of the
\(\{\neg,\to\}\)-fragment,
\[
 \Gamma\cup\{\varphi\}\derives_{\mathsf H_{\mathsf{CPL}}}\psi
 \quad\Longleftrightarrow\quad
 \Gamma\derives_{\mathsf H_{\mathsf{CPL}}}\varphi\to\psi ,
\]
where \(\derives_{\mathsf H_{\mathsf{CPL}}}\) is derivability by finite
sequences in this calculus.  This is Theorem~\ref{thm:p02-hilbert-deduction};
its proof, in Annex~\ref{annex:p2.calc.2}, is an induction on the given
derivation that uses only A1, A2 and MP, and from right to left it is one MP
step.  The theorem holds for this premise-sensitive formula-sequence calculus
because MP is its only rule: uniform substitution instantiates axiom
schemas and does not license arbitrary substitution into an assumption.

\paragraph{Translations.}
Hilbert derivations on this sheet are turned into derivations on
Sheet~\ref{sheet:nd-cpl} and into propositional derivations on
Sheet~\ref{sheet:lk}, and conversely through the translation of
Definition~\ref{def:p02-hilbert-translation}, by
Theorems~\ref{thm:p02-w02b-p02-010-hilbert-to-nd},
\ref{thm:p02-w02b-p02-010-nd-to-hilbert},
\ref{thm:p02-w02b-p02-010-hilbert-to-lk} and
\ref{thm:p02-w02b-p02-010-lk-to-hilbert}; the resulting derivability
equivalence is Corollary~\ref{thm:p02-w02b-p02-010-hilbert-equivalence}.
